% !TEX root = ../main.tex
\corrigesChapitre{25}{Photons et effet photoélectrique}
\corrigesSouples

\begin{blocLivre}{cor-25-1}
\begin{corrige}{1}{L'énergie d'un photon}
\demarche{$E=\dfrac{h\,c}{\lambda}$, puis division par \qty{1,60e-19}{J} pour obtenir des électronvolts.}
\begin{questions}
\item Violet : \res{\qty{3,11}{eV}} ; rouge : \res{\qty{1,78}{eV}}.
\item UV : \res{\qty{4,89}{eV}} ; infrarouge : \res{\qty{1,24}{eV}}.
\item Il faut $E\geqslant\qty{3,6}{eV}$, soit $\lambda\leqslant\dfrac{h\,c}{E}=\res{\qty{345}{nm}}$ : seuls les photons \res{\text{ultraviolets}} peuvent
rompre cette liaison.
\item Un photon ultraviolet a assez d'énergie pour \res{\text{rompre des liaisons}} des molécules biologiques (ADN) ; un
photon infrarouge n'excite que l'agitation thermique.
\end{questions}
\end{corrige}
\end{blocLivre}
\begin{blocLivre}{cor-25-2}
\begin{corrige}{2}{Les seuils photoélectriques}
\demarche{$\lambda_0=\dfrac{h\,c}{W_\mathrm{ext}}$, avec $W_\mathrm{ext}$ converti en joules.}
\begin{questions}
\item C'est l'énergie \res{\text{minimale}} qu'il faut fournir pour extraire un électron du métal.
\item Césium : \res{\qty{592}{nm}} ; sodium : \res{\qty{545}{nm}} ; zinc : \res{\qty{289}{nm}} ; cuivre : \res{\qty{264}{nm}}.
\item Le \res{\text{césium}} (de 400 à \qty{592}{nm}) et le \res{\text{sodium}} (de 400 à \qty{545}{nm}), pour une partie du visible.
\item Son seuil, le plus élevé en longueur d'onde, permet de détecter presque toute la lumière visible.
\end{questions}
\end{corrige}
\end{blocLivre}
\begin{blocLivre}{cor-25-3}
\begin{corrige}{3}{Des électrons arrachés au sodium}
\demarche{$E_\mathrm{c,max}=h\nu-W_\mathrm{ext}$, puis $v=\sqrt{2E_\mathrm{c}/m}$.}
\begin{questions}
\item $E=\dfrac{hc}{\lambda}=\res{\qty{3,11}{eV}}$.
\item $h\nu=W_\mathrm{ext}+E_\mathrm{c,max}$ : $E_\mathrm{c,max}=3{,}11-2{,}28=\res{\qty{0,83}{eV}}=\res{\qty{1,32e-19}{J}}$.
\item $v=\sqrt{\dfrac{2\times\num{1.32e-19}}{\num{9.11e-31}}}=\res{\qty{5,4e5}{m.s^{-1}}}$.
\item À \qty{600}{nm}, $E=\qty{2,07}{eV}<W_\mathrm{ext}$ : \res{\text{aucun électron}} n'est émis, même si l'on double l'intensité, qui ne
fait que doubler le nombre de photons, sans changer leur énergie.
\end{questions}
\end{corrige}
\end{blocLivre}
\begin{blocLivre}{cor-25-4}
\begin{corrige}{4}{Onde ou particule ?}
\demarche{Modèle ondulatoire : énergie continue et cumulable ; modèle particulaire : énergie par paquets $h\nu$.}
\begin{questions}
\item Le métal finirait par recevoir assez d'énergie : des électrons seraient émis \res{\text{au bout d'un certain temps}}.
\item En dessous de la fréquence seuil, \res{\text{aucun électron}} n'est jamais émis.
\item Chaque électron absorbe un seul photon ; si $h\nu<W_\mathrm{ext}$, aucun photon n'a \res{\text{assez d'énergie}} pour l'extraire.
\item Les deux : la lumière a une \res{\text{double nature}}. Les interférences illustrent l'aspect ondulatoire, l'effet
photoélectrique l'aspect particulaire.
\end{questions}
\end{corrige}
\end{blocLivre}
\begin{blocLivre}{cor-25-5}
\begin{corrige}{5}{Combien de photons ?}
\demarche{Nombre de photons par seconde $=$ puissance / énergie d'un photon.}
\begin{questions}
\item $E=\dfrac{\num{6.63e-34}\times\num{3.00e8}}{\num{532e-9}}=\res{\qty{3,74e-19}{J}}$.
\item $N=\dfrac{\num{1.0e-3}}{\num{3.74e-19}}=\res{\num{2.7e15}}$ photons par seconde.
\item $N=\dfrac{1000}{\num{3.62e-19}}\approx\res{\num{2.8e21}}$ photons par seconde.
\item Les photons sont si nombreux que la lumière paraît \res{\text{continue}}.
\end{questions}
\end{corrige}
\end{blocLivre}
\begin{blocLivre}{cor-25-6}
\begin{corrige}{6}{À la manière de Millikan}
\demarche{$e\,U_\mathrm{a}=h\nu-W_\mathrm{ext}$ : droite de pente $h/e$.}
\begin{questions}
\item $E_\mathrm{c,max}=e\,U_\mathrm{a}=h\nu-W_\mathrm{ext}$, d'où $\res{U_\mathrm{a}=\dfrac{h}{e}\,\nu-\dfrac{W_\mathrm{ext}}{e}}$.
\item $a=\dfrac{1{,}69-0{,}57}{(9{,}6-6{,}9)\times10^{14}}=\res{\qty{4,15e-15}{V.s}}$ et $b=0{,}57-a\times\num{6.9e14}=\res{\qty{-2,29}{V}}$.
\item $h=a\,e=\res{\qty{6,64e-34}{J.s}}$ ; $W_\mathrm{ext}=-b\,e$, soit \res{\qty{2,29}{eV}}.
\item $\nu_0=\dfrac{-b}{a}=\res{\qty{5,5e14}{Hz}}$. Les mesures, alignées sur une droite de pente $h/e$, confirment
\res{\text{quantitativement}} la relation d'Einstein et donnent une valeur de $h$ conforme aux autres mesures.
\end{questions}
\end{corrige}
\end{blocLivre}
\begin{blocLivre}{cor-25-7}
\begin{corrige}{7}{L'électroscope de Hallwachs}
\demarche{Seuls les photons d'énergie supérieure à \qty{4,3}{eV} (UV, $\lambda<\qty{289}{nm}$) arrachent des électrons au zinc.}
\begin{questions}
\item Les photons ultraviolets arrachent des électrons à la plaque, qui \res{\text{perd sa charge négative}} : les feuilles
retombent.
\item Les photons visibles ont une énergie \res{\text{inférieure}} à \qty{4,3}{eV} : aucun électron n'est arraché, quelle que soit
la puissance.
\item \res{\text{Non}} : les électrons arrachés sont aussitôt rappelés par la plaque positive ; sa charge ne change pas.
\item Le verre arrête les ultraviolets : la plaque \res{\text{ne se décharge plus}}.
\end{questions}
\end{corrige}
\end{blocLivre}
\begin{blocLivre}{cor-25-8}
\begin{corrige}{8}{Le rendement d'un panneau en Python}
\demarche{$\eta=\dfrac{P_\mathrm{max}}{E\times S}$.}
\begin{questions}
\item \code{P} contient les puissances électriques $UI$ ; \code{k} est l'indice de la mesure de \res{\text{puissance maximale}}.
\item La puissance lumineuse reçue vaut $E\times S=1000\times0{,}025=\qty{25}{W}$ : le rendement, en pourcentage, est
\res{\dfrac{100\,P_\mathrm{max}}{25}}.
\item On multiplie les intensités par 0,8 dans la liste \code{I} et l'on remplace \code{1000} par \code{800} : la puissance
électrique et la puissance reçue diminuent toutes deux de 20~\%, le rendement \res{\text{ne change pas}} (13,7~\%).
\item Le point de fonctionnement dépend de l'appareil branché : il ne coïncide pas en général avec le point de puissance
maximale. Les onduleurs ajustent en permanence la tension pour \res{\text{suivre ce maximum}}.
\end{questions}
\end{corrige}
\end{blocLivre}
\begin{blocLivre}{cor-25-9}
\begin{corrige}{9}{La diode électroluminescente}
\demarche{Émission : photon d'énergie $hc/\lambda$ ; nombre de photons $=$ puissance lumineuse / énergie d'un photon.}
\begin{questions}
\item Par \res{\text{émission}} : $E=\dfrac{hc}{\lambda}=\res{\qty{1,97}{eV}}$.
\item Tension de seuil voisine de \res{\qty{2,0}{V}}.
\item $N=\dfrac{0{,}35}{\num{3.16e-19}}=\res{\num{1.1e18}}$ photons par seconde.
\item $\eta=\dfrac{0{,}35}{1{,}0}=\res{\qty{35}{\percent}}$ ; le reste est converti en \res{\text{énergie interne}} : la LED chauffe.
\end{questions}
\end{corrige}
\end{blocLivre}
\begin{blocLivre}{cor-25-10}
\begin{corrige}{10}{Les limites d'une cellule au silicium}
\demarche{Absorption si $\dfrac{hc}{\lambda}\geqslant\qty{1,12}{eV}$ ; énergie en excès $=\dfrac{hc}{\lambda}-\qty{1,12}{eV}$.}
\begin{questions}
\item $\lambda_\mathrm{max}=\dfrac{hc}{\qty{1,12}{eV}}=\res{\qty{1,11}{\micro m}}$.
\item $E=\qty{0,83}{eV}<\qty{1,12}{eV}$ : il \res{\text{n'est pas absorbé}}.
\item Excès : $3{,}11-1{,}12=\res{\qty{1,99}{eV}}$, soit \res{\qty{64}{\percent}} de son énergie, perdue.
\item Les photons de grande longueur d'onde ne sont pas absorbés et ceux de courte longueur d'onde perdent leur énergie
en excès : une grande partie de l'énergie solaire est \res{\text{inexploitable}}, quelle que soit la qualité de la cellule.
\end{questions}
\end{corrige}
\end{blocLivre}
\begin{blocLivre}{cor-25-11}
\begin{corrige}{11}{Le capteur d'un appareil photo}
\demarche{Énergie $=$ éclairement $\times$ surface $\times$ durée ; nombre de photons $=$ énergie / $(hc/\lambda)$.}
\begin{questions}
\item $E=0{,}10\times(\num{4e-6})^2\times0{,}010=\res{\qty{1,6e-14}{J}}$.
\item $N=\dfrac{\num{1.6e-14}}{\num{3.62e-19}}=\res{\num{4.4e4}}$ photons, soit $0{,}60\,N=\res{\num{2.7e4}}$ électrons.
\item \res{\text{Non}}, de peu : la saturation serait atteinte pour une pose de $0{,}010\times\dfrac{\num{30000}}{\num{26600}}\approx\res{\qty{11}{ms}}$.
\item L'\res{\text{effet photoélectrique}} (interne) : chaque photon absorbé libère un électron.
\end{questions}
\end{corrige}
\end{blocLivre}
\begin{blocLivre}{cor-25-12}
\begin{corrige}{12}{Des panneaux solaires sur le toit}
\demarche{Rendement $=$ puissance crête / (\qty{1000}{W.m^{-2}} $\times$ surface) ; énergie annuelle $=\eta\times$ énergie reçue $\times0{,}8$.}
\begin{questions}
\item $\eta=\dfrac{400}{1000\times1{,}7}=\res{\qty{23,5}{\percent}}$.
\item $10\times\qty{400}{W}=\res{\qty{4,0}{kW}}$ crête.
\item $E=0{,}235\times1300\times17\times0{,}8\approx\res{\qty{4,2e3}{kWh}}$ : un peu moins que les \qty{4500}{kWh} consommés, la
production couvre \res{\text{plus de 90~\%}} des besoins sur l'année.
\item L'éclairement dépend de la hauteur du Soleil, donc de l'heure et de la saison, et de la météo ; la nuit, la
production est \res{\text{nulle}}.
\end{questions}
\end{corrige}
\end{blocLivre}
\begin{blocLivre}{cor-25-13}
\begin{corrige}{13}{Détecter un photon unique}
\demarche{$\lambda_0=hc/W_\mathrm{ext}$ ; $E_\mathrm{c,max}=hc/\lambda-W_\mathrm{ext}$ ; charge $=10^6\times e$.}
\begin{questions}
\item $\lambda_0=\dfrac{hc}{\qty{2,0}{eV}}=\res{\qty{622}{nm}}$.
\item $E=\qty{2,96}{eV}$, d'où $E_\mathrm{c,max}=\res{\qty{0,96}{eV}}$.
\item $Q=10^6\times\num{1.60e-19}=\res{\qty{1,6e-13}{C}}$ ; on compte en moyenne \res{250} impulsions pour \num{1000} photons.
\item L'énergie d'un seul photon est infime ; l'effet photoélectrique la convertit en un électron, que l'appareil
\res{\text{multiplie}} jusqu'à un signal mesurable.
\end{questions}
\end{corrige}
\end{blocLivre}
\begin{blocLivre}{cor-25-14}
\begin{corrige}{14}{La cellule photoélectrique au césium}
\demarche{Photons par seconde $=P/(hc/\lambda)$ ; $I=$ (électrons par seconde) $\times\,e$.}
\begin{questions}
\item $N=\dfrac{\num{1.0e-6}}{\num{4.42e-19}}=\res{\num{2.26e12}}$ photons par seconde.
\item $I=0{,}20\times\num{2.26e12}\times\num{1.60e-19}=\res{\qty{72}{nA}}$.
\item $E_\mathrm{c,max}=2{,}76-2{,}1=\res{\qty{0,66}{eV}}$.
\item Puissance doublée : $I$ \res{\text{double}} (\qty{145}{nA}), $E_\mathrm{c,max}$ ne change pas. Lumière rouge : $E=\qty{1,78}{eV}<\qty{2,1}{eV}$,
\res{\text{plus aucun courant}}.
\end{questions}
\end{corrige}
\end{blocLivre}
\begin{blocLivre}{cor-25-15}
\begin{corrige}{15}{La couleur du laser hélium-néon}
\demarche{Énergie du photon $=$ écart entre les niveaux, puis $\lambda=hc/E$.}
\begin{questions}
\item \res{\text{Émission}} : $E=20{,}66-18{,}70=\res{\qty{1,96}{eV}}$.
\item $\lambda=\dfrac{hc}{E}=\res{\qty{634}{nm}}$ : lumière \res{\text{rouge}}.
\item $\lambda=\dfrac{1239{,}84}{1{,}96}=\qty{632,6}{nm}$, très proche de la valeur tabulée : l'écart initial venait des
\res{\text{arrondis}} de $h$, $c$ et $e$.
\item L'émission correspond à une seule transition entre deux niveaux précis : une seule longueur d'onde. Une lampe à
incandescence émet un \res{\text{spectre continu}}, d'origine thermique.
\end{questions}
\end{corrige}
\end{blocLivre}
